%% kapa: exam=1 round=24 year=2013 period=2 form=A subject=appraisal_law subject_ko=부동산관계법규 seq=29 number=69
%% source: Q-Net 감정평가사 기출문제 (cst00309) artlSeq=5073802 file=extracted/1008/5073802/2013년 제24회 감정평가사 1차시험 기출문제/2교시(A형).hwp
%% conversion: kordoc 4.24.51 -> tools/gen1.py (auto); review=proofread-sonnet; flags=none
\begin{question}{69}
건축법령상 조정(調停) 및 재정(裁定)에 관한 설명으로 옳지 않은 것은?
\begin{choices}
  \item 재정은 문서로써 하여야 한다.
  \item 조정은 3명의 위원으로 구성되는 조정위원회에서 한다.
  \item 건축분쟁전문위원회는 재정신청이 된 사건을 조정에 회부하는 것이 적합하다고 인정하면 직권으로 직접 조정할 수 있다.
  \item 당사자가 재정에 불복하여 소송을 제기한 경우 시효의 중단과 제소기간의 산정에 있어서는 재정신청을 재판상의 청구로 본다.
  \item 재정위원회의 회의는 구성원 과반수의 출석으로 열고 출석한 위원 과반수의 찬성으로 의결한다.
\end{choices}
\end{question}
\answer{5}
